\chapter{IDENTIFIKASI MASALAH DAN IDE INOVATIF}

\section{Identifikasi Masalah}

\subsection{Kompleksitas Deployment pada Kubernetes}

Pengelolaan aplikasi pada klaster Kubernetes memerlukan penanganan berbagai komponen yang saling terkait, termasuk konfigurasi jaringan, penyimpanan, dan keamanan. Proses deployment manual sering kali rentan terhadap kesalahan konfigurasi dan sulit untuk dilacak.

\subsection{Manajemen Rahasia yang Tidak Aman}

Rahasia seperti kredensial database dan token API sering kali disimpan dalam bentuk plaintext atau di luar kontrol versi, sehingga sulit untuk diaudit dan rentan terhadap kebocoran data.

\subsection{Kurangnya Otomatisasi dalam Deployment}

Tanpa mekanisme otomatisasi, setiap perubahan pada aplikasi atau infrastruktur memerlukan intervensi manual dari operator, yang meningkatkan waktu deployment dan risiko kesalahan manusia.

\section{Pengembangan Ide Inovatif}

\subsection{GitOps sebagai Solusi}

GitOps menawarkan pendekatan deklaratif di mana seluruh status sistem dideskripsikan dalam file konfigurasi yang tersimpan di repositori Git. Pendekatan ini memungkinkan:

\begin{itemize}
    \item Audit trail lengkap untuk setiap perubahan.
    \item Kemudahan rollback ke versi sebelumnya.
    \item Kolaborasi yang lebih baik melalui alur kerja pull request.
\end{itemize}

\subsection{Inovasi yang Diusulkan}

Produk yang dikembangkan dalam penelitian ini mengintegrasikan beberapa teknologi untuk membentuk solusi GitOps yang komprehensif:

\begin{enumerate}
    \item \textbf{ArgoCD} untuk orkestrasi deployment otomatis dengan dukungan sync wave.
    \item \textbf{Sealed Secrets} untuk enkripsi rahasia yang aman untuk disimpan di Git.
    \item \textbf{Cloudflare Tunnel} untuk akses eksternal tanpa membuka port inbound.
    \item \textbf{NetworkPolicies} untuk isolasi jaringan berbasis default-deny.
\end{enumerate}
